% --------------------------------------------------
\subsubsection{Тогтмол вариацын зорилтот матриц}
% --------------------------------------------------

Тогтмол вариацын зорилтот матриц нь бүх хөрөнгийн вариацыг нэг нийтлэг утгаар төлөөлүүлж, хөрөнгүүдийн хоорондын ковариацыг тэг гэж үздэг. Уг зорилтот матрицыг
\[
T_{\mathrm{var}}
=
\mu I
\]
гэж тодорхойлно. Энд $I$ нь $N\times N$ хэмжээст нэгж матриц, харин
\[
\mu
=
\frac{1}{N}\sum_{i=1}^{N}s_{ii}
=
\frac{\operatorname{tr}(S)}{N}
\]
нь хөрөнгүүдийн өгөөжийн түүврийн вариацуудын дундаж юм. Иймээс тогтмол вариацын зорилтот матриц
\[
T_{\mathrm{var}}
=
\begin{pmatrix}
	\mu & 0 & \cdots & 0\\
	0 & \mu & \cdots & 0\\
	\vdots & \vdots & \ddots & \vdots\\
	0 & 0 & \cdots & \mu
\end{pmatrix}
\]
хэлбэртэй байна \citep{ledoitwolf2004a}.

Тогтмол вариацын зорилтот матрицад суурилсан шугаман агшилтын үнэлэгчийг
\[
\widehat{\Sigma}_{\mathrm{var}}
=
(1-\widehat{\alpha}_{\mathrm{var}})S
+
\widehat{\alpha}_{\mathrm{var}}T_{\mathrm{var}},
\qquad
0\leq\widehat{\alpha}_{\mathrm{var}}\leq1,
\]
гэж тодорхойлно. Энд $\widehat{\alpha}_{\mathrm{var}}$ нь тогтмол вариацын зорилтот матрицад харгалзах агшилтын эрчим юм.

\paragraph{Агшилтын эрчмийн үнэлгээ.}

Ledoit--Wolf-ийн аргаар агшилтын эрчмийг үнэлэхдээ түүврийн ковариацын матриц болон зорилтот матрицын хоорондын зай, мөн түүврийн ковариацын матрицын элементүүдийн түүвэрлэлтийн хэлбэлзлийг ашиглана \citep{ledoitwolf2004a}.

Эхлээд түүврийн ковариацын матриц $S$ болон тогтмол вариацын зорилтот матриц $T_{\mathrm{var}}$ хоорондын квадрат Фробениусын зайг
\[
\widehat{\gamma}_{\mathrm{var}}
=
\left\|
S-T_{\mathrm{var}}
\right\|_F^2
\]
гэж тодорхойлно. Уг зайг матрицын элементүүдээр илэрхийлбэл
\[
\widehat{\gamma}_{\mathrm{var}}
=
\sum_{i=1}^{N}
\sum_{j=1}^{N}
(s_{ij}-t_{ij})^2,
\]
энд $s_{ij}$ болон $t_{ij}$ нь харгалзан $S$ болон $T_{\mathrm{var}}$ матрицын $(i,j)$-р элементүүд юм.

Дараа нь түүврийн ковариацын матрицын элементүүдийн түүвэрлэлтийн хэлбэлзлийг тооцохын тулд
\[
z_{tij}
=
(y_{ti}-\overline{Y}_i)
(y_{tj}-\overline{Y}_j)
\]
гэж тэмдэглэе. Энд $y_{ti}-\overline{Y}_i$ нь $t$-р ажиглалт дахь $i$-р хөрөнгийн өгөөж болон тухайн хөрөнгийн дундаж өгөөжийн зөрүүг илэрхийлнэ. Тэгвэл түүврийн ковариацын элемент
\[
s_{ij}
=
\frac{1}{n}
\sum_{t=1}^{n}z_{tij}
\]
байна.

Түүврийн ковариацын матрицын элементүүдийн түүвэрлэлтийн хэлбэлзлийг нэгтгэн хэмжих хэмжигдэхүүнийг
\[
\widehat{\pi}
=
\frac{1}{n}
\sum_{t=1}^{n}
\sum_{i=1}^{N}
\sum_{j=1}^{N}
(z_{tij}-s_{ij})^2
\]
гэж тооцно. Энд $z_{tij}-s_{ij}$ нь $t$-р ажиглалтаар тооцогдох $i$ болон $j$-р хөрөнгийн өгөөжийн хамтын өөрчлөлт нь тэдгээрийн түүврийн ковариац $s_{ij}$-аас хэр зэрэг зөрж байгааг илэрхийлнэ.

Тогтмол вариацын зорилтот матрицад харгалзах агшилтын эрчмийг
\[
\widehat{\alpha}_{\mathrm{var}}
=
\max
\left\{
0,
\min
\left(
1,
\frac{\widehat{\pi}}
{n\widehat{\gamma}_{\mathrm{var}}}
\right)
\right\}
\]
гэж тооцно \citep{ledoitwolf2004a}. Ийнхүү $0\leq\widehat{\alpha}_{\mathrm{var}}\leq1$ нөхцөл хангагдана.

Эцэст нь тогтмол вариацын зорилтот матрицад суурилсан агшилтын ковариацын матрицыг
\[
\widehat{\Sigma}_{\mathrm{var}}
=
(1-\widehat{\alpha}_{\mathrm{var}})S
+
\widehat{\alpha}_{\mathrm{var}}T_{\mathrm{var}}
\]
гэж олно.

\paragraph{Тооцоолох алгоритм.}

Тогтмол вариацын зорилтот матрицад суурилсан агшилтын ковариацын матрицыг дараах үе шаттайгаар тооцно.

\begin{enumerate}
	\item Хөрөнгө тус бүрийн дундаж өгөөжийг тооцно.
	\[
	\overline{Y}_i
	=
	\frac{1}{n}
	\sum_{t=1}^{n}y_{ti},
	\qquad
	i=1,2,\ldots,N.
	\]
	
	\item Түүврийн ковариацын матрицыг байгуулна.
	\[
	s_{ij}
	=
	\frac{1}{n}
	\sum_{t=1}^{n}
	(y_{ti}-\overline{Y}_i)
	(y_{tj}-\overline{Y}_j),
	\]
	томьёогоор $S=(s_{ij})$ матрицыг тооцно.
	
	\item Түүврийн вариацуудын дундажийг тооцно.
	\[
	\mu
	=
	\frac{1}{N}
	\sum_{i=1}^{N}s_{ii}
	=
	\frac{\operatorname{tr}(S)}{N}.
	\]
	
	\item Тогтмол вариацын зорилтот матрицыг байгуулна.
	\[
	T_{\mathrm{var}}
	=
	\mu I.
	\]
	
	\item Түүврийн ковариацын матриц болон зорилтот матрицын хоорондын квадрат Фробениусын зайг тооцно.
	\[
	\widehat{\gamma}_{\mathrm{var}}
	=
	\left\|
	S-T_{\mathrm{var}}
	\right\|_F^2.
	\]
	
	\item Түүврийн ковариацын матрицын элементүүдийн түүвэрлэлтийн хэлбэлзлийг тооцно. Эхлээд
	\[
	z_{tij}
	=
	(y_{ti}-\overline{Y}_i)
	(y_{tj}-\overline{Y}_j)
	\]
	гэж тооцоод,
	\[
	\widehat{\pi}
	=
	\frac{1}{n}
	\sum_{t=1}^{n}
	\sum_{i=1}^{N}
	\sum_{j=1}^{N}
	(z_{tij}-s_{ij})^2
	\]
	гэж олно.
	
	\item Агшилтын эрчмийг тооцно.
	\[
	\widehat{\alpha}_{\mathrm{var}}
	=
	\max
	\left\{
	0,
	\min
	\left(
	1,
	\frac{\widehat{\pi}}
	{n\widehat{\gamma}_{\mathrm{var}}}
	\right)
	\right\}.
	\]
	
	\item Агшилтын ковариацын матрицыг байгуулна.
	\[
	\widehat{\Sigma}_{\mathrm{var}}
	=
	(1-\widehat{\alpha}_{\mathrm{var}})S
	+
	\widehat{\alpha}_{\mathrm{var}}T_{\mathrm{var}}.
	\]
\end{enumerate}